\section{STARKs: FRI, Polynomial Commitments, and AIR for Neural Networks}

\subsection{Overview}

STARKs differ fundamentally from circuit-based proof systems.  
Instead of compiling computation into static constraints, they model it as a dynamic process over time using the \emph{Algebraic Intermediate Representation (AIR)}.

\medskip

\noindent
\textbf{Pipeline.}
\[
\text{execution trace}
\;\longrightarrow\;
\text{polynomial representation}
\;\longrightarrow\;
\text{constraint composition}
\;\longrightarrow\;
\text{FRI low-degree proof}.
\]

\subsection{From Trace to Polynomials}

Let the execution trace have $W$ columns:
\[
T_1, T_2, \dots, T_W.
\]

Each column is interpolated into a polynomial:
\[
f_j(X), \qquad \deg(f_j) \le n-1,
\]
such that for a multiplicative subgroup $G=\{g^0,\dots,g^{n-1}\}$,
\[
f_j(g^t) = T_j[t].
\]

\subsection{Constraint Polynomials}

A transition rule such as:
\[
x_{t+1} = x_t^2 + 1
\]
becomes:
\[
f_x(gX) - f_x(X)^2 - 1 = 0.
\]

Boundary constraints are encoded via divisibility:
\[
f_x(1) = x_0
\quad \Longleftrightarrow \quad
(X-1) \mid (f_x(X) - x_0).
\]

\subsection{Quotient Construction}

If a constraint polynomial $C(X)$ must vanish on a set $S$, then:
\[
C(X) = Z_S(X) \cdot Q(X),
\quad
Z_S(X) = \prod_{a \in S} (X-a).
\]

Thus:
\[
Q(X) = \frac{C(X)}{Z_S(X)}
\]
must be a low-degree polynomial.

Multiple constraints are combined:
\[
C_{\mathrm{comp}}(X) =
\sum_i \alpha_i \frac{C_i(X)}{Z_i(X)}.
\]

\subsection{FRI as a Polynomial Commitment}

The prover evaluates polynomials over an extended domain and commits via Merkle trees.

\medskip

\noindent
\textbf{Key step: folding.}

Given:
\[
p_0(X) = g_1(X^2) + X g_2(X^2),
\]
define:
\[
p_1(X) = g_1(X) + \beta g_2(X).
\]

This reduces degree and domain size. Repeating yields:
\[
p_0 \to p_1 \to p_2 \to \cdots.
\]

\medskip

\noindent
\textbf{Interpretation.}  
FRI proves that the committed evaluations correspond to a low-degree polynomial.

\subsection{AIR for Convolutional Neural Networks}

Consider a $2\times2$ convolution applied to a $3\times3$ input.

\paragraph{Trace design.}

Each row corresponds to one output location.

Columns:
\[
p_1,p_2,p_3,p_4,\quad
k_1,k_2,k_3,k_4,\quad
s,\quad y.
\]

\paragraph{Local constraint.}

\[
s_t = \sum_{i=1}^4 k_i p_{i,t} + b.
\]

\paragraph{Activation (example).}

\[
y_t = s_t^2.
\]

\paragraph{Transition constraint.}

The patch shifts across the image:
\[
\text{patch}_{t+1} = \text{Shift}(\text{patch}_t).
\]

\paragraph{Boundary constraints.}

\begin{itemize}
\item first row = first patch
\item kernel values fixed
\item last row = final output location
\end{itemize}

\medskip

\noindent
\textbf{Observation.}  
CNNs are well-suited to AIR because of repeated local structure.

\subsection{AIR for Transformer Attention}

Consider one attention head:
\[
Q = XW_Q,\quad K = XW_K,\quad V = XW_V.
\]

\subsubsection{Stage 1: Linear projections}

\[
Q_{ir} = \sum_j X_{ij} (W_Q)_{jr}.
\]

Implemented via accumulation constraints over rows.

\subsubsection{Stage 2: Attention scores}

\[
s_{ij} = \sum_r Q_{ir} K_{jr}.
\]

Each $(i,j)$ pair is computed through a sequence of rows.

\subsubsection{Stage 3: Normalization}

Approximate softmax:
\[
u_{ij} = \widetilde{\exp}(s_{ij}),\quad
z_i = \sum_j u_{ij},\quad
a_{ij} z_i = u_{ij}.
\]

\subsubsection{Stage 4: Output}

\[
O_{ir} = \sum_j a_{ij} V_{jr}.
\]

\subsection{Comparison: CNN vs Transformer in AIR}

\begin{center}
\begin{tabular}{l|c|c}
Feature & CNN & Transformer \\
\hline
Interaction & local & global \\
Trace size & moderate & large \\
Nonlinearity & simple & complex (softmax) \\
AIR suitability & high & moderate
\end{tabular}
\end{center}

\subsection{Key Insight}

\begin{quote}
AIR represents computation as a discrete dynamical system:
\[
s_{t+1} = F(s_t),
\]
and correctness reduces to verifying low-degree polynomial structure.
\end{quote}

\subsection{Discussion}

\begin{itemize}
\item FRI provides a transparent polynomial commitment via low-degree testing.
\item CNNs align naturally with AIR due to locality.
\item Transformers are expressible but expensive due to global interactions.
\end{itemize}

\medskip

\noindent
\textbf{Conclusion.}  
STARKs combine AIR and FRI to enable scalable, transparent proofs of complex computations, including machine learning models.